\paragraph{Native preparation of the same ground and its charged response.}
The original cube ground and its local charged spectral measure now have
an explicit native source preparation and readout.  The exact
highest-weight ground $G$, with $N=4$ and $E_G=-\sqrt3/4$, has 62 nonzero
hidden-up occupation amplitudes $(a+b\sqrt3)/36$, where $a,b$ are integers
and $ab=0$.  Its squared amplitudes are therefore exact integers divided
by 1296.  These source-derived weights and signs determine an eight-bit
binary amplitude tree; no fitted spectrum or unknown-state oracle is used.
The existing compiler emits 5,379 literal native pulses.  An independent
rational interval check consumes every event, inverts the Walsh layers,
and verifies the sine and cosine of each conditional rotation against
the exact signed square-root ratios.  Disjoint control blocks give a
maximum error per level, and telescoping the levels and complete native
load/export maps bounds the full initialized preparation, including aborts.

Eight blank down modes are included by an exact isometric embedding into
all 65,536 Fock rows.  A further 184 native pulses apply the inverse
onsite frame.  Each local frame word is checked on its full four-state
Fock space, including vacuum and doublon phases, using
$e^{is\pi P/4}=(I+isP)/\sqrt2$ and Gaussian-integer products.
Its literal-angle and native-instrument errors are separately enclosed.
The resulting complete preparation-state error is below $10^{-10}$;
the retained exact bound is approximately $8.67\times10^{-13}$.

At the origin the hidden and original frames coincide, and the down
mode is empty in $G$.  Thus
\[
\gamma_1=Z_0X_1=c_1+c_1^\dagger,
\qquad \gamma_1G=c_1^\dagger G,
\qquad \|\gamma_1G\|=1.
\]
The Jordan--Wigner sign is essential: replacing $Z_0X_1$ by $X_1$
produces an orthogonal probe because the exact up occupation is $1/2$.
This unitary creation preserves the entire initial exterior CAR state:
disjoint even observables commute, and exterior odd expectations vanish
in both definite-parity states.  The origin-down occupation changes from
zero to one and total number from four to five.  Its mean added energy
is $1/2$ with variance $67/256$; this is a localized initial defect with a
conserved number flag, not a persistent localization or mass-eigenstate
claim.
A 27-pulse pointer word prepares
$(|0\rangle G+|1\rangle\gamma_1G)/\sqrt2$.
Its controlled Majorana is the product of commuting Pauli rotations in
$\exp[i\pi(I-Z_{\rm pointer})(I-\gamma_1)/4]$.
This explicitly supplied native qubit service maintains a phase relation
between the $N=4$ and $N=5$ arms; it does not derive a physical
fermion-parity reference from passive repair.

Both arms evolve under the same original Hamiltonian.  With a second
controlled Majorana and pointer analyzer of phase $\phi$,
\begin{align*}
C(t)&=\langle\gamma_1G,e^{-it(H-E_G)}\gamma_1G\rangle,\\
P_0(\phi,t)&=\tfrac12\left[1+\operatorname{Re}
                         (e^{i\phi}C(t))\right].
\end{align*}
The ground-energy phase is supplied by the evolving reference arm;
no separate scalar offset is inserted into the source law.
The two executed settings have $t=1$ and $\phi=0,\pi/2$.
Each setting uses two half-time segments, each with 32 symmetric
splitting steps and 337,920 native pulses, followed by a 29-pulse
analyzer.  Including its ground preparation and pointer, a complete
setting uses 681,459 pulses on the full 17-bit field/pointer register.

\paragraph{Protection and complete output accounting.}
The five C repair seams follow the first time segment, and the five E
seams follow the second.  The complete repair theorem acts as identity
on $N4\oplus N5$, including cross-number coherences: all native branches
vanish there, and the outside branch is identity.  This is passive
protection during the waits.  Preparation and control words are not
claimed to remain in this subspace at every intermediate pulse.
The final inverse Majorana can create an $N=6$ component after the
charged particle has moved; no repair is applied after that analyzer.

Every source first-export failure is retained as its own terminal
record.  The numerical execution also retains all 130 repair outcome
slots, including off-subspace outputs from control error, and all six
final pointer outcomes.  A native repair outcome terminates with that
record; continuation uses the outside branch without renormalization.
The lossless archive retains the prepared full ground, both full final
states and every first-failure weight table.  Individual failure vectors
remain addressable from their actual source prefixes.

\paragraph{Connection to the resolved spectrum.}
All field and pointer words are consumed against the exact source
templates.  Their complete operator errors, the preparation error and
the full-Fock Strang errors of both time segments give a complete output
error below $3/5000$ for either setting.  Contractivity through the
complete repair and readout channels preserves the telescoping bound.
The same bound controls each unconditioned pointer probability.
Independently, the 78 certified spectral intervals evaluate
\[
C(t)=\sum_{j=1}^{78}w_j e^{-i\Delta E_jt}
\]
with rational sine and cosine enclosures.  Adding the source error gives
certified pointer intervals, checked against the literal execution and
an independent exponential of the original Hamiltonian.  The
$\pi/2$ outcome is separated from $1/2$; the $U=0$ comparator has the
purely real response $\cos(\sqrt3t/16)$.  Zero-time native runs check the
analyzers as well.  These two settings measure the charged characteristic
function.  They do not independently resolve all 78 lines, select an
autonomous ground, prepare a thermodynamic vacuum, establish laboratory
units, or turn finite energy gaps into relativistic particle masses.
